% !TEX root = ../aa_SoM_Chapters.tex

%% HARRIS: Chapter 10

% ö = \%c3\%b6
% ä = \%C3\%A4

\xdef\sectiontitle{\IIsectiontitle} %aiheuttama
\TOCrefs

%%%%%%%%%%%%%%%
\begin{frame}[label=2_5_a]%[label=2_4_TOC1]
\frametitle{\texorpdfstring{\culine[\sectiontitleulinecolor]{\sectiontitle}}{\sectiontitle} }

%\vspace{0.2cm}
\vspace{\chapterTOCvksip}
\begin{itemize}[leftmargin=0.5cm,itemsep=3mm]
    \item[2.1] \hyperlink{2_1_a}{\IIaaSubsectiontitle}
    \item[2.2] \hyperlink{2_2_a}{\IIaSubsectiontitle}
    \item[2.3] \hyperlink{2_3_a}{\IIbSubsectiontitle}	
    \item[2.4] \hyperlink{2_4_a}{\IIcSubsectiontitle}	
    \item[2.5] \hyperlink{2_5_a}{\CDAlert[red]{\IIdSubsectiontitle}}
    \item[2.6] \hyperlink{2_6_a}{\IIeSubsectiontitle}
    \item[2.7] \hyperlink{2_7_a}{\IIfSubsectiontitle}
    \item[2.8] \hyperlink{2_8_a}{\IIgSubsectiontitle}
    \item[2.9] \hyperlink{2_9_a}{\IIhSubsectiontitle}
    \item[2.10] \hyperlink{2_10_a}{\IIiSubsectiontitle}
\end{itemize}

%\endofslide 
\end{frame}
%%%%%%%%%%%%%%%

\xdef\sectiontitle{\IIdSubsectiontitle} 

%%%%%%%%%%%%%%%%
%\begin{frame}
%\frametitle{\texorpdfstring{\culine[\sectiontitleulinecolor]{\sectiontitle}}{\sectiontitle} }
%\framesubtitle{Energy bands}
%
%\vspace{1.2mm}
%\begin{minipage}{9.5cm}
%\begin{itemize}[itemsep=2.5mm] \small
%	% 6.5 cm = midway, 10 cm = 3/4 
%
%	\item The atoms in solid material have many valence electrons\appear{, but only a fraction of them can contribute to electrical conduction}
%
%\item Also, in electrical insulators\appear{, there needs to be} \appear{(must be)} \appear{a mechanism to prohibit free movement of the valence electrons}
%
%\item Clearly, the potential well describing the atom that contains the electrons\appear{, is not a simple well}\appear{, but it somehow restricts the electrons to \CDAlert[red]{energy bands}}
%
%\item With two wells close to each other\appear{, the energy levels split forming a \CDAlert[blue]{fine structure}} 
%
%\item For $N$ atoms close together\appear{, we get energy bands of $N$ levels, respectively} 
%
%\item This is a good starting point for understanding a crystal with positive ions and electrons attracted by them
%\implies{ Mode detailed and satisfactory treatment follows \appear{in \\
%		the course on solid-state physics}}
%\end{itemize}
%\end{minipage}
%
%\xdef\loc{north east}
%\begin{tikzpicture}[remember picture,overlay] % anchor=south
%    \node (A) [yshift=\figYshift,xshift=\figXshift, anchor=\loc] at (current page.\loc) {\includegraphics[page=1,width=4cm]{Figures_booklet/SOM_10_Bonding_figures/SOM_Bonding_Figure_1.png}}; %scale=\figscale. % 8cm max height
%\end{tikzpicture}
%
%\xdef\loc{south east}
%\begin{tikzpicture}[remember picture,overlay] % anchor=south
%    \node (A) [yshift=-0.25*\figYshift,xshift=\figXshift, anchor=\loc] at (current page.\loc) {\includegraphics[page=1,width=4cm]{Figures_booklet/SOM_10_Bonding_figures/SOM_Bonding_Figure_2.png}}; %scale=\figscale. % 8cm max height
%\end{tikzpicture}
%
%\endofslide 
%%\endofsection
%\end{frame}
%%%%%%%%%%%%%%%%

%%%%%%%%%%%%%%%
\begin{frame}
\frametitle{\texorpdfstring{\culine[\sectiontitleulinecolor]{\sectiontitle}}{\sectiontitle} }
\framesubtitle{Energy bands}

\vspace{1.2mm}
%\begin{minipage}{9.5cm}
\begin{itemize}[itemsep=1.8mm] \small
	% 6.5 cm = midway, 10 cm = 3/4 
%	\item Atoms in solids have many valence electrons\appear{, but only a fraction of them can contribute to electrical conduction}

\item Only a small fraction of the valence electrons in solids contribute to electrical conduction
\item In dielectrics\appear{, some mechanism must prohibit free movement of valence electrons}

\item Clearly, the potential well describing the atom that contains the electrons\appear{, is not a simple well}\appear{, but it somehow restricts the electrons to \CDAlert[red]{energy bands}}

\item With two wells close to each other\appear{, the energy levels split forming a \CDAlert[blue]{fine structure}} 

\item For $N$ atoms close together\appear{, we get energy bands of $N$ levels, respectively} 

\item This is a good starting point for understanding a crystal with positive ions and electrons attracted by them
\implies{\small Mode detailed and satisfactory treatment follows \appear{in \\
		the course on solid-state physics}}
\end{itemize}
%\end{minipage}

%\xdef\loc{north east}
%\begin{tikzpicture}[remember picture,overlay] % anchor=south
%    \node (A) [yshift=\figYshift,xshift=\figXshift, anchor=\loc] at (current page.\loc) {\includegraphics[page=1,width=4cm]{Figures_booklet/SOM_10_Bonding_figures/SOM_Bonding_Figure_1.png}}; %scale=\figscale. % 8cm max height
%\end{tikzpicture}

\xdef\loc{south west}
\begin{tikzpicture}[remember picture,overlay] % anchor=south
    \node (A) [yshift=-0.15*\figYshift,xshift=-\figXshift, anchor=\loc] at (current page.\loc) {\includegraphics[page=1,width=8.5cm]{Figures_booklet/SOM_10_Bonding_figures/SOM_Bonding_Figure_1.png}}; %scale=\figscale. % 8cm max height
\end{tikzpicture}

\xdef\loc{south east}
\begin{tikzpicture}[remember picture,overlay] % anchor=south
    \node (A) [yshift=-0.25*\figYshift,xshift=\figXshift, anchor=\loc] at (current page.\loc) {\includegraphics[page=1,width=4cm]{Figures_booklet/SOM_10_Bonding_figures/SOM_Bonding_Figure_2.png}}; %scale=\figscale. % 8cm max height
\end{tikzpicture}

\endofslide 
%\endofsection
\end{frame}
%%%%%%%%%%%%%%%

%%%%%%%%%%%%%%%
\begin{frame}
\frametitle{\texorpdfstring{\culine[\sectiontitleulinecolor]{\sectiontitle}}{\sectiontitle} }

\begin{minipage}{9cm}
\begin{itemize}
	\item Let's first compare energies in infinite 1D quantum wells of sizes $4 L$:
$$\semismall
\begin{aligned} 
\appear{E_{4,4L}}  &\equal \appear{16} \frac{\pi^{2} \hbar^{2}}{2 m(4 L)^{2}} \appear{=}\frac{16}{16} \frac{\pi^{2} \hbar^{2}}{2 m L^{2}} \\
\appear{E_{3,4 L}} &\equal \appear{9} \frac{\pi^{2} \hbar^{2}}{2 m(4 L)^{2}} \equal \frac{9}{16} \frac{\pi^2 \hbar^{2}}{2 m L^{2}} \\
\appear{E_{2,4 L}} &\equal \appear{4} \frac{\pi^{2} \hbar^{2}}{2 m(4 L)^{2}}  \equal \frac{4}{16} \frac{\pi^{2} \hbar^{2}}{2 m L^{2}} \\
\appear{E_{1,4 L}} &\equal \frac{\pi^{2} \hbar^{2}}{2 m(4 L)^{2}} \equal \frac{1}{16} \frac{\pi^{2} \hbar^{2}}{2 m L^{2}}
\end{aligned}
$$
\appear{and  for a well of width $L$:}\vspace{3mm}
\[ \semismall
%\begin{aligned}
\appear{E_{1, L}} \equal \frac{\pi^{2} \hbar^{2}}{2 m L^{2}}\appear{\,,} \qquad
\appear{E_{2, L}} \equal \appear{4} \frac{\pi^{2} \hbar^{2}}{2 m L^{2}} \equal \appear{4 E_{1, L}}\appear{\,,} \qquad
\appear{E_{3, L}} \equal  \appear{9 E_{1, L}}%\appear{9} \frac{\pi^{2} \hbar^{2}}{2 m L^{2}}  \equal
%\end{aligned}
\]

\end{itemize}
\end{minipage}

\xdef\loc{south east}
\begin{tikzpicture}[remember picture,overlay] % anchor=south
    \node (A) [yshift=0*\figYshift,xshift=\figXshift, anchor=\loc] at (current page.\loc) {\includegraphics[page=1,height=9cm]{Figures_booklet/SOM_10_Bonding_figures/SOM_Bonding_Figure_23.png}}; %scale=\figscale. % 8cm max height
\end{tikzpicture}

\endofslide 
%\endofsection
\end{frame}
%%%%%%%%%%%%%%%

%%%%%%%%%%%%%%%
\begin{frame}
\frametitle{\texorpdfstring{\culine[\sectiontitleulinecolor]{\sectiontitle}}{\sectiontitle} }
 
%\vspace{2.5mm}
\begin{minipage}{9cm}
\begin{itemize}
	\item Now we have finite wells, that can be populated with electrons from a crystal consisting of 4 atoms\appear{, e.g. the 12$^{th}$ energy state is given by}
\[
E_{12,4 L} \equal \frac{144}{16} \frac{\pi^{2} \hbar^{2}}{2 m L^{2}}  \equal \appear{9} \frac{\pi^{2} \hbar^{2}}{2 m L^{2}}
\]

\item Figure 23 shows the principle of what happens with a 4-atom crystal (in the middle)

\item The bands are formed, and their energies cluster around single-atom energies

\item Note that energies are smaller when placing 12 electrons into an well of width $4L$\appear{, than placing the electrons into several wells of width $L$}
\end{itemize}
\end{minipage}

\xdef\loc{south east}
\begin{tikzpicture}[remember picture,overlay] % anchor=south
    \node (A) [yshift=0*\figYshift,xshift=\figXshift, anchor=\loc] at (current page.\loc) {\includegraphics[page=1,height=9cm]{Figures_booklet/SOM_10_Bonding_figures/SOM_Bonding_Figure_23.png}}; %scale=\figscale. % 8cm max height
\end{tikzpicture}

\endofslide 
%\endofsection
\end{frame}
%%%%%%%%%%%%%%%

%%%%%%%%%%%%%%%
\begin{frame}
\frametitle{\texorpdfstring{\culine[\sectiontitleulinecolor]{\sectiontitle}}{\sectiontitle} }

\begin{minipage}{9cm}
\begin{itemize}
	\item The states naturally divide into 4 levels: 
	
	\item[] \begin{itemize}[itemsep=2.2mm] \semismall
	\item For example, the highest wave function at the band $n=2$\appear{, has 8 antinodes}\appear{, and is zero exactly in between the atoms (/positive ions), where the periodic potential energy is at maximum}
	\item Also, it does not differ much from the single atom wave function
	
	\item The lowest wave function of the $n=2$ band\appear{, however, has more or less 5 antinodes}\appear{, and the maxima of the wave function are located at the interatomic points} 
	
	\item And, it has the energy much higher than what it would be without the periodic variation (in a large well)
	\end{itemize}

\end{itemize}
\end{minipage}

\xdef\loc{south east}
\begin{tikzpicture}[remember picture,overlay] % anchor=south
    \node (A) [yshift=0*\figYshift,xshift=\figXshift, anchor=\loc] at (current page.\loc) {\includegraphics[page=1,height=9cm]{Figures_booklet/SOM_10_Bonding_figures/SOM_Bonding_Figure_23.png}}; %scale=\figscale. % 8cm max height
\end{tikzpicture}

\endofslide 
%\endofsection
\end{frame}
%%%%%%%%%%%%%%%

%%%%%%%%%%%%%%%
\begin{frame}
\frametitle{\texorpdfstring{\culine[\sectiontitleulinecolor]{\sectiontitle}}{\sectiontitle} }

\begin{minipage}{9cm}
\begin{itemize}
	\item So far, we have seen that electrons in the 4-atom crystal\appear{, made of 4 finite quantum wells close to each other}\appear{, have \CDAlert[red]{energy bands} set by the energy levels of the \CDAlert[red]{L-well = a single atom}}

\item However, the \CDAlert[fuchsia]{wave functions} with their wavelength\appear{, resemble more the wave functions in the \CDAlert[fuchsia]{4L-well = well size for 4 atoms}}

\item Next, we analyze the wavelengths of the electron wave functions \appear{by using the concept of wave number as given by}
\[
k  \equal \frac{2 \pi}{\lambda}
\]
\end{itemize}
\end{minipage}

\xdef\loc{south east}
\begin{tikzpicture}[remember picture,overlay] % anchor=south
    \node (A) [yshift=0*\figYshift,xshift=\figXshift, anchor=\loc] at (current page.\loc) {\includegraphics[page=1,height=9cm]{Figures_booklet/SOM_10_Bonding_figures/SOM_Bonding_Figure_23.png}}; %scale=\figscale. % 8cm max height
\end{tikzpicture}

\endofslide 
%\endofsection
\end{frame}
%%%%%%%%%%%%%%%

%%%%%%%%%%%%%%%
\begin{frame}
\frametitle{\texorpdfstring{\culine[\sectiontitleulinecolor]{\sectiontitle}}{\sectiontitle} }

\begin{itemize}
	\item Having obtained the wave number $k$\appear{, we can calculate the kinetic energy of the electrons in each state:} 
\end{itemize}

\begin{minipage}{7.3cm}
\begin{itemize}
\item[] \[
E_{kin} \equal \frac{p^{2}}{2 m} \equal \frac{\hbar^{2} k^{2}}{2 m}
\]

\item[] \begin{itemize}[itemsep=2.5mm] \small
 \item We now plot the \CDAlert[fuchsia]{actual energy} versus $k$\appear{, along with \CDAlert[black]{kinetic energy}, curve is a \CDAlert[black]{parabola}} \appear{(}$\appear{E_{kin}} \propto \appear{k^{2}} $\appear{)} \appear{and is marked using a \CDAlert[black]{dashed line}} 
	 \item We see that the energy at each level \CDAlert[red]{exceeds} \CDAlert[red]{the kinetic energy}\appear{, the difference being potential energy of each state} $ \appear{\textrm{((}E_{\text {tot}}=}\appear{E_{\text {kin}}} \plus \appear{E_{\text {pot}} \textrm{))}}$ 
	 
	 \item The lowest state in each band has significantly \CDAlert[tunired]{high potential energy}
\end{itemize}

\end{itemize}
\end{minipage}

\xdef\loc{south east}
\begin{tikzpicture}[remember picture,overlay] % anchor=south
    \node (A) [yshift=-0.25*\figYshift,xshift=\figXshift, anchor=\loc] at (current page.\loc) {\includegraphics[page=1,height=7.6cm]{Figures_booklet/SOM_10_Bonding_figures/SOM_Bonding_Figure_24.png}}; %scale=\figscale. % 8cm max height
\end{tikzpicture}

\endofslide 
%\endofsection
\end{frame}
%%%%%%%%%%%%%%%

%%%%%%%%%%%%%%%
\begin{frame}
\frametitle{\texorpdfstring{\culine[\sectiontitleulinecolor]{\sectiontitle}}{\sectiontitle} }

\begin{itemize}
	\item We now increase the number of atoms in the crystal\appear{, atomic spacing remaining as $a$} 
	\implies{ \CDAlert[red]{Energy bands are formed} }  
\end{itemize}

\vspace{1.3mm}
\begin{minipage}{7.3cm}
\begin{itemize}
	
	\item[] \begin{itemize}[itemsep=2.1mm] \small
	 \item For the state at the top of each band\appear{, the nodes of wave function hit the midpoint between the atoms $(\mathrm{x})$}

\item Each state\appear{/wavefunction} \appear{in band $n$ holds $n$ antinodes}\appear{, and the spacing between atoms is $\mathrm{n}$ times half wavelength}
\[
a  \equal \appear{n \cdot} \frac{1}{2} \appear{\lambda} \quad \Rightarrow \quad \appear{k} \equal \appear{n} \frac{\pi}{a}
\]

\item At the bottom of each band\appear{, the lowest state has high potential energy} 

%\item $N$ being large, the wavelengths at top and bottom of energy bands are practically the same

\end{itemize}

\end{itemize}
\end{minipage}

\xdef\loc{south east}
\begin{tikzpicture}[remember picture,overlay] % anchor=south
    \node (A) [yshift=-0.25*\figYshift,xshift=\figXshift, anchor=\loc] at (current page.\loc) {\includegraphics[page=1,height=8.2cm]{Figures_booklet/SOM_10_Bonding_figures/SOM_Bonding_Figure_25.png}}; %scale=\figscale. % 8cm max height
\end{tikzpicture}

\endofslide 
%\endofsection
\end{frame}
%%%%%%%%%%%%%%%

%%%%%%%%%%%%%%%%
%\begin{frame}
%\frametitle{\texorpdfstring{\culine[\sectiontitleulinecolor]{\sectiontitle}}{\sectiontitle} }
%
%\begin{itemize}
%	\item Obviously: The jump between $n$ and $n+1$ states occurs sharply at 
%\end{itemize}
%
%	% 6.5 cm = midway, 10 cm = 3/4 
%\vspace{2.5mm}
%\begin{minipage}{7cm}
%\begin{itemize}
%	\item[] $k=n \Frac{\pi}{a}$. \appear{The top of band $n$ has $n$ antinodes per atom}\appear{, and there are $N$ atoms overall}\appear{, then this state has a total of $n N$ antinodes}
%
%\item The top of band $n+1$ has $n+1$ antinodes per atom\appear{, and since there are $N$ atoms overall}\appear{, then this state has a total of $(n+1) N$ antinodes.} \appear{So the number of states from one band top to the next band top to will be $N$}
%
%\item So clearly: \Alert[red]{Each band consists of N} 
%		\CDAlert[red]{states}
%\end{itemize}
%\end{minipage}
%
%\xdef\loc{south east}
%\begin{tikzpicture}[remember picture,overlay] % anchor=south
%    \node (A) [yshift=-0.25*\figYshift,xshift=\figXshift, anchor=\loc] at (current page.\loc) {\includegraphics[page=1,height=8.2cm]{Figures_booklet/SOM_10_Bonding_figures/SOM_Bonding_Figure_25.png}}; %scale=\figscale. % 8cm max height
%\end{tikzpicture}
%
%\endofslide 
%%\endofsection
%\end{frame}
%%%%%%%%%%%%%%%%

%%%%%%%%%%%%%%%
\begin{frame}
\frametitle{\texorpdfstring{\culine[\sectiontitleulinecolor]{\sectiontitle}}{\sectiontitle} }

\begin{itemize}
	\item Since $N$ is large, the \CDAlert[blue]{energy band is practically a continuum}

\item When $a$ is not too small\appear{, the energy band spreads above and below the energies of a single atom system}

\item The extent of \CDAlert[red]{spreading of the band depends on} how much \CDAlert[red]{overlap} there is \CDAlert[red]{between the wave functions} \appear{of adjacent ($1^{\text{st}}$ neighbouring) atoms}
\end{itemize}
% 6.5 cm = midway, 10 cm = 3/4 
\vspace{2.5mm}
\begin{minipage}{4cm}
\begin{itemize}
	\item The \CDAlert[tunired]{band width}: 
\item[] \begin{itemize}
	\item does not depend on the number of atoms 
	\item depends on the \CDAlert[tunired]{atomic spacing $a$}
\end{itemize}

\end{itemize}
\end{minipage}


\xdef\loc{south  east}
\begin{tikzpicture}[remember picture,overlay] % anchor=south
    \node (A) [yshift=-0.25*\figYshift,xshift=\figXshift, anchor=\loc] at (current page.\loc) {\includegraphics[page=1,width=9.75cm]{Figures_booklet/SOM_10_Bonding_figures/SOM_Bonding_Figure_26.png}}; %scale=\figscale. % 8cm max height
\end{tikzpicture}

\endofslide 
%\endofsection
\end{frame}
%%%%%%%%%%%%%%%

%%%%%%%%%%%%%%%
\begin{frame}
\frametitle{\texorpdfstring{\culine[\sectiontitleulinecolor]{\sectiontitle}}{\sectiontitle} }
%\framesubtitle{\texorpdfstring{\culine[\sectiontitleulinecolor]{Summary}}{Summary} }
\framesubtitle{\texorpdfstring{\DUTline[\sectiontitleulinecolor]{Summary}}{Summary} }
	
%\vspace{2.5mm}
\begin{minipage}{8.2cm} %,itemindent=0pt,leftmargin=*
\begin{itemize}
	\item[] \begin{itemize}[itemsep=1.5mm,leftmargin=0mm] \small
	 \item The crystal is macroscopic ($N$ large)\appear{, therefore a continuum of energies arise} 
	 
	 \item Microscopic ionic potential energies divide the continuum into $n$ energy bands\appear{, related to individual atomic states}
	
	\item The band structure \appear{($E$ vs. wave number $k$)} \appear{is often presented by a graph called \CDAlert[fuchsia]{dispersion relation}} \appear{(Fig.~27)}  
	
	\item Note that $k$ can be positive or negative\appear{, depending on the direction of the movement of the state}
	
	\item The (\Alert[fuchsia]{fuchsia}) graph \CDAlert[red]{almost follows} the (dashed black) parabola
	\vspace{-1mm}
	 \implies{ ~Electron has \CDAlert[red]{mostly kinetic energy},\\
	  i.\,e. \CDAlert[tunired]{can move freely} }
	 %in the macroscopic infinite well  
	\item Exceptions of this arise near $\appear{k} \equal \appear{n} \frac{\pi}{a}$

\item Final note, so far we have only discussed the states\appear{, not how they will be filled by electrons}
\end{itemize}

\end{itemize}
\end{minipage}


\xdef\loc{south east}
\begin{tikzpicture}[remember picture,overlay] % anchor=south
    \node (A) [yshift=-0.25*\figYshift,xshift=\figXshift, anchor=\loc] at (current page.\loc) {\includegraphics[page=1,height=8cm]{Figures_booklet/SOM_10_Bonding_figures/SOM_Bonding_Figure_27.png}}; %scale=\figscale. % 8cm max height
\end{tikzpicture}

\endofslide 
%\endofsection
\end{frame}
%%%%%%%%%%%%%%%

%%%%%%%%%%%%%%%
\begin{frame}
\frametitle{\texorpdfstring{\culine[\sectiontitleulinecolor]{\sectiontitle}}{\sectiontitle} }
\framesubtitle{\texorpdfstring{\DUTline[\sectiontitleulinecolor]{Electrical conduction}}{Electrical conduction} }

\begin{itemize}
	\item \Alert[red]{Classical} model ('Drude model') of electrical conduction:\\
	\item[] \begin{itemize}[itemsep=1.8mm] \small
	 \item In 1900, three years after \CDAlert[red]{the discovery of electrons by Thomson} 
	 \item Paul Drude proposed that there would be \CDAlert[blue]{free electrons} in metals
	 \appear{\xdef\loc{south east}%
\begin{tikzpicture}[remember picture,overlay]% anchor=south
    \node (A) [yshift=-0.25*\figYshift,xshift=\figXshift, anchor=\loc] at (current page.\loc) {\includegraphics[page=1,width=5cm]{Figures_booklet/SOM_10_Bonding_figures/Tikz_SOM_Bonding_e_conduction_schematic_fixed}};%scale=\figscale. % 8cm max height
\end{tikzpicture}}
	 \item Model acts as a basis for more sophisticated models \appear{explaining electric and thermal conduction} \appear{(\CDAlert[red]{Drude–Lorentz model})}\appear{, and can even be made compatible with quantum mechanics} \appear{(\CDAlert[violet]{Drude–Sommerfeld model})}
\end{itemize}


\item This model, developed further by Lorentz \appear{(Drude–Lorentz model)}\appear{, is called classical model of electrical conduction}
\end{itemize}

\vspace{2.5mm}
\begin{minipage}{8.5cm}
\begin{itemize}

\item  In the absence of electrical field\appear{, the mean speed of conduction electrons is given by:} 
% 6.5 cm = midway, 10 cm = 3/4 
\[
< v >\,  \equal \appear{\textrm{((}} \frac{8 k_{B} T}{\pi m_{e}} \appear{\textrm{))}^{1/2}}  \equal  \appear{1.08 \times} \appear{10^{5}} \appear{\mathrm{~m} / \mathrm{s}} \appear{\,, \quad T=293~K}
\]
\end{itemize}
\end{minipage}
\vspace{2.5mm}


\endofslide 
%\endofsection
\end{frame}
%%%%%%%%%%%%%%%


%%%%%%%%%%%%%%%
\begin{frame}
\frametitle{\texorpdfstring{\culine[\sectiontitleulinecolor]{\sectiontitle}}{\sectiontitle} }
\framesubtitle{\texorpdfstring{\DUTline[\sectiontitleulinecolor]{Electrical conduction}}{Electrical conduction} }

\begin{itemize}[itemsep=1.8mm]
	\item[] 
$$
< v >\,= \textrm{((}\Frac{8 k_{B} T}{\pi m_{e}} \textrm{))}^{1/2}=1.08 \times 10^{5} \mathrm{~m} / \mathrm{s} \,, \quad T=293~K
$$

\item But due to collisions within the material\appear{, the trajectory of the electron is random}

\item A potential difference across the ends of a conductor produces an electrical field inside the conductor
\implies{ This exerts a (Lorentz) force on the electron $F=e E$ } 

\item Free electrons move freely within the material\appear{, \\
	but collide with positive ions} 
\item The electron will be drifting due to \appear{(against)} \\
	\appear{the applied electric field} \appear{(see the analogy with\\
	the ball rolling in spiky inclined plane)}\appear{, and \\
	are accelerated between each collision}
\end{itemize}

\xdef\loc{south east}%
\begin{tikzpicture}[remember picture,overlay]% anchor=south
    \node (A) [yshift=-0.25*\figYshift,xshift=\figXshift, anchor=\loc] at (current page.\loc) {\includegraphics[page=1,width=5cm]{Figures_booklet/SOM_10_Bonding_figures/Tikz_SOM_Bonding_e_conduction_schematic_fixed}};%scale=\figscale. % 8cm max height
\end{tikzpicture}

\endofslide 
%\endofsection
\end{frame}
%%%%%%%%%%%%%%%

%%%%%%%%%%%%%%%
\begin{frame}
\frametitle{\texorpdfstring{\culine[\sectiontitleulinecolor]{\sectiontitle}}{\sectiontitle} }

\begin{itemize}
	\item The average time between subsequent collisions is $\tau$
	\item Within this time\appear{, the electron has gained velocity of} $\appear{\mathrm{e} \mathrm{E}}  \appear{\tau} \appear{/ \mathrm{m}_{\mathrm{e}}}$

\item When there are many collisions\appear{, there will be an average \CDAlert[red]{drift velocity}} \appear{which will eventually cause a current density $\boldsymbol{j}$}

\item Basic electrodynamics gives:
$$ \seminormal
%\begin{aligned}
\appear{j} \equal \frac{\text {charge}}{\text{time} \cdot \text{area}} \equal \frac{\text{charge}}{\text {distance} \cdot \text{area}} \appear{\cdot} \frac{\text{distance}}{\text {time}}
 \equal \frac{e \times \text {charge number}}{\text {volume}} \appear{\cdot} \frac{\text{distance}}{\text {time}}
%\end{aligned}
$$
\appear{Thus: }
\[
 j \equal \appear{e} \appear{\rho} \appear{v_{\text {drift}}}  \equal \appear{e \rho} \frac{e E}{m_{e}} \appear{\tau}  \equal \frac{e^{2} \rho \tau}{m_{e}} \appear{E}
\]
\item[] where $\rho$ is the number of free charge carriers per unit volume
\end{itemize}

\endofslide 
%\endofsection
\end{frame}
%%%%%%%%%%%%%%%

%%%%%%%%%%%%%%%
\begin{frame}
\frametitle{\texorpdfstring{\culine[\sectiontitleulinecolor]{\sectiontitle}}{\sectiontitle} }

\begin{itemize}
	\item The equation above gives, \appear{in practice, \CDAlert[red]{Ohm's law}}\appear{, where the \Alert[blue]{proportionality constant is the electrical conductivity of the material}}\appear{, and its reciprocal would be resistivity of the material:} 
\[
j  \equal \appear{\sigma} \appear{E\,,} \qquad \appear{ \text {where}} \qquad \appear{\sigma}  \equal \frac{e^{2} \appear{\rho \tau}}{m_{e}}
\]

\item This Drude model does not give correct values for the conductivity

\item In addition, the values for drift velocity will be unrealistically small \appear{(at least one order of magnitude smaller than experimentally estimated values) }

\end{itemize}

%\xdef\loc{south east}
%\begin{tikzpicture}[remember picture,overlay] % anchor=south
%    \node (A) [yshift=-0.25*\figYshift,xshift=\figXshift, anchor=\loc] at (current page.\loc) {\includegraphics[page=4,width=5cm]{Figures_booklet/Tikz_SOM_Bonding_figures_fixed}}; %scale=\figscale. % 8cm max height
%\end{tikzpicture}

\endofslide 
%\endofsection
\end{frame}
%%%%%%%%%%%%%%%

%%%%%%%%%%%%%%%
\begin{frame}
\frametitle{\texorpdfstring{\culine[\sectiontitleulinecolor]{\sectiontitle}}{\sectiontitle} }

\begin{itemize}
	
\item But the \Alert[red]{classical model is compatible with quantum mechanics}: 
\item[] \begin{itemize}
	 \item \CDAlert[tunired]{Electrical resistance is seen to originate from collisions} 
	 
	 \item But quantum mechanically\appear{, the \CDAlert[blue]{collision time} is much larger than that obtained classically with electrons colliding with every positive ion} 
	 
	 \item In fact, the electrons do not simply 'bump off' the ions 
	 \item The periodic crystal allows free electrons to act almost like free particles 
	 \item Also, \CDAlert[fuchsia]{electrons are waves}\appear{, so the positive\\
	 	 ions have a little effect}\appear{, thus the \Alert[red]{collisions need to be considered differently}}
\end{itemize}
 \implies{ Need to take into account the \CDAlert[fuchsia]{band structure of the solid} }
\end{itemize}

%\xdef\loc{south east}
%\begin{tikzpicture}[remember picture,overlay] % anchor=south
%    \node (A) [yshift=-0.25*\figYshift,xshift=\figXshift, anchor=\loc] at (current page.\loc) {\includegraphics[page=4,width=5cm]{Figures_booklet/Tikz_SOM_Bonding_figures}}; %scale=\figscale. % 8cm max height
%\end{tikzpicture}


\endofslide 
%\endofsection
\end{frame}
%%%%%%%%%%%%%%%

%%%%%%%%%%%%%%%
\begin{frame}
\frametitle{\texorpdfstring{\culine[\sectiontitleulinecolor]{\sectiontitle}}{\sectiontitle} }

\begin{itemize}[itemsep=1.7mm]
\item Valence electrons are assumed to be freely moving in a periodic quantum well, in which case their energies are those of fermion gas

	\item In metals the fermionic electrons\appear{ fill the energy bands from lowest states}\appear{, following the \CDAlert[red]{Pauli exclusion principle}}

\end{itemize}

% 6.5 cm = midway, 10 cm = 3/4 
\vspace{1.7mm}
\begin{minipage}{3.5cm}
\begin{itemize}[itemsep=1.7mm]

\item \appear{Energy of highest occupied state is called \CDAlert[fuchsia]{Fermi energy $E_{F}$}}

\item In a conductor\appear{, the highest band is only partially filled}\appear{, and this is the band where the conduction electrons reside}

%	\item If no external electric field is applied, the random net movement of the electrons would be zero. If an external electric field is applied, then some of the electrons shift oppositely in field's direction, and slightly above $E_{F}$
\end{itemize}
\end{minipage}


\xdef\loc{south east}
\begin{tikzpicture}[remember picture,overlay] % anchor=south
    \node (A) [yshift=-0.25*\figYshift,xshift=\figXshift, anchor=\loc] at (current page.\loc) {\includegraphics[page=1,width=10cm]{Figures_booklet/SOM_10_Bonding_figures/SOM_Bonding_Figure_28.png}}; %scale=\figscale. % 8cm max height
\end{tikzpicture}

\endofslide 
%\endofsection
\end{frame}
%%%%%%%%%%%%%%%

%%%%%%%%%%%%%%%
\begin{frame}
\frametitle{\texorpdfstring{\culine[\sectiontitleulinecolor]{\sectiontitle}}{\sectiontitle} }

\begin{itemize}[itemsep=1.7mm]
%	\item In metals the electrons, as fermions, fill the energy bands from lowest states, following the \CDAlert[red]{Pauli exclusion principle}
%
%\item In a conductor, the topmost band is only partially filled, and this is the band where the conduction electrons reside
\item Valence electrons are assumed to be freely moving in a periodic quantum well, in which case their energies are those of fermion gas

	\item In metals the fermionic electrons\appear{ fill the energy bands from lowest states}\appear{, following the \CDAlert[red]{Pauli exclusion principle}}

\end{itemize}

% 6.5 cm = midway, 10 cm = 3/4 
\vspace{1.7mm}
\begin{minipage}{3.5cm}
\begin{itemize}[itemsep=1.7mm]
	\item If no external electric field is applied\appear{, the random net movement of the electrons is zero}
	\item Once an external electric field is applied\appear{, some of the electrons shift accordingly moving slightly above $E_{F}$}
\end{itemize}
\end{minipage}


\xdef\loc{south east}
\begin{tikzpicture}[remember picture,overlay] % anchor=south
    \node (A) [yshift=-0.25*\figYshift,xshift=\figXshift, anchor=\loc] at (current page.\loc) {\includegraphics[page=1,width=10cm]{Figures_booklet/SOM_10_Bonding_figures/SOM_Bonding_Figure_28.png}}; %scale=\figscale. % 8cm max height
\end{tikzpicture}

%\endofslide 
\endofsection
\end{frame}
%%%%%%%%%%%%%%%